% Generated draft: requires mathematical and source review.
\documentclass[11pt]{article}
\usepackage[margin=25mm]{geometry}
\usepackage{fontspec}
\setmainfont{Latin Modern Roman}
\usepackage{amsmath,amssymb,mathtools,unicode-math}
\setmathfont{Latin Modern Math}
\usepackage{enumitem,needspace,xurl,hyperref}
\setlength{\parindent}{0pt}
\setlength{\parskip}{5pt}
\setlength{\emergencystretch}{3em}
\newcommand{\N}{\mathbb{N}}
\newcommand{\Z}{\mathbb{Z}}
\newcommand{\Q}{\mathbb{Q}}
\newcommand{\R}{\mathbb{R}}
\newcommand{\C}{\mathbb{C}}
\newcommand{\F}{\mathbb{F}}
\newcommand{\A}{\mathbb{A}}
\newcommand{\PP}{\mathbb{P}}
\newcommand{\E}{\mathbb{E}}
\newcommand{\eps}{\varepsilon}
\newcommand{\dd}{\,\mathrm{d}}
\begin{document}
\section*{UnsolvedMath ID: 30004339}
\textbf{Source problem code:} catalogue entry
\subsection{Definitions and mathematical statement}
Regard a finitely generated group \(\Gamma\) as a discrete group. For an essentially free probability-measure-preserving action \(\Gamma\curvearrowright(X,\mu)\) on a standard nonatomic probability space, let \(R_\Gamma\) denote its orbit equivalence relation. A graphing of \(R_\Gamma\) is a countable family \(\Phi=(\varphi_i:A_i\to B_i)_i\) of measure-preserving partial isomorphisms whose graphs generate \(R_\Gamma\), modulo null sets. Define
\[
\operatorname{cost}(R_\Gamma)=\inf_{\Phi}\sum_i\mu(A_i).
\]
One standard convention for group cost is
\[
\operatorname{cost}_{\inf}(\Gamma)=\inf_{\Gamma\curvearrowright(X,\mu)}\operatorname{cost}(R_\Gamma),
\]
where the infimum ranges over all such essentially free actions. Let \(\mathcal N\Gamma\) be the group von Neumann algebra with its canonical trace. The first \(L^2\)-Betti number is
\[
\beta_1^{(2)}(\Gamma)=\dim_{\mathcal N\Gamma}H_1\!\left(\mathcal N\Gamma\otimes_{\Z\Gamma}C_*(E\Gamma;\Z)\right),
\]
where \(E\Gamma\) is a contractible free \(\Gamma\)-CW complex and \(\dim_{\mathcal N\Gamma}\) is von Neumann dimension.

\textbf{Supplied question, conditional on the group-cost convention above.} Is it true that
\[
\forall\,\Gamma\text{ finitely generated},\qquad
\operatorname{cost}_{\inf}(\Gamma)=1+\beta_1^{(2)}(\Gamma)?
\]
\textbf{Unresolved qualifications.} The record does not define group cost or specify whether the intended assertion concerns its infimum over actions or the cost of every essentially free probability-measure-preserving action. These quantifiers must be distinguished. Moreover, the supplied quantifier includes finite groups. Under the convention above, a finite group has cost \(1-|\Gamma|^{-1}\) and first \(L^2\)-Betti number zero, so the displayed unrestricted formulation is incompatible with that boundary case. An infinitude hypothesis or an altered formula would change the supplied statement; neither is silently inserted here. The intended formulation therefore requires source clarification.
\subsection{Short English statement}
The supplied question asks whether every finitely generated group has cost equal to one plus its first \(L^2\)-Betti number. Its intended meaning requires clarification of the cost convention, the quantifier over actions, and the inclusion of finite groups.
\subsection{Sources}
\begin{itemize}
\item Supplied local UnsolvedMath registry record, permanent ID \(30004339\), title ``Cost and First \(L^2\) Betti Numbers,'' source problem code \texttt{catalogue entry}, in the Oberwolfach Reports Open Problems collection. The recorded proposed year is 2019; no individual proposer is supplied.
\item Supplied report citation: \emph{Groups, Dynamics, and Approximation} (2020), Oberwolfach Reports, DOI 10.4171/owr/2019/55, pp. 3409--3458. \url{https://doi.org/10.4171/owr/2019/55}. The background describes these reports as records of mathematical workshops and their open-problem discussions.
\item Supplied literature reference: Gaboriau--Seward, ``cost and square-summable Betti numbers.'' \url{https://mathweb.ucsd.edu/~bseward/Files/costbettientropy.pdf}. This is the supplied descriptive label; full bibliographic details are not provided.
\item Supplied literature reference: ``OWR report on groups, dynamics, and approximation.'' \url{https://ems.press/content/serial-article-files/46835}.
\item Literature assessment date: \texttt{2026-08-22}. This is supplied dated curation derived from the OWR record, not an independently verified assessment. It labels the problem open and states that the general cost equality remains open, while cost one is known for Kazhdan groups and other special classes. The scope qualifications of these assertions, including infinitude where necessary, require verification. The supplied background explicitly requests independent verification of time-sensitive status and cited claims.
\item The registry labels the statement \texttt{corrected\_verified}, with the explicit qualification that only formula delimiters, notation, grammar, and typography were repaired, preserving mathematical content and scope. This label does not establish mathematical correctness or resolve the formulation issues above.
\end{itemize}
\end{document}
